\DocumentMetadata{pdfversion=1.7,pdfstandard=a-2b,lang=en-US}
\RequirePackage[l2tabu, orthodox]{nag}
\documentclass[10pt, a4paper]{article}
\input{../header.tex}
\input{../headerSuppressMetaPDF.tex}
%\pdfusetitle
\usepackage{fancyvrb}

\hypersetup{
  pdfinfo={
    % in newest version also in custom dictionary 
    %Author      ={Ernst Reißner},
    %Title       ={The DVI-format and the program DVItype},
    Subject     ={DVI and DVItype},
    Keywords    ={LaTeX;dvi;DVItype},
    unicode % TBD: clarify 
  }
}
%\pdfcatalog{/Lang (de-DE)}




\newcommand{\bs}{\symbol{'134}} % backslash in typewriter T1/T1 

\IfPackageLoadedTF{tex4ht}{
  \newcommand{\cmd}[2][]{\texttt{\bs{}#2#1}}
  % \newcommand{\pkg}[1]{\texttt{#1}}% without indexing 
  % \newcommand{\tool}[1]{\texttt{#1}}% without indexing 
  % \newcommand{\param}[1]{\texttt{#1}}% without indexing 
}
{
  \newcommand{\cmd}[2][]{\texttt{\bs{}#2#1}\sindex[cmd]{#2@\texttt{\bs{}#2}}}
  % \newcommand{\pkg}[1]{\texttt{#1}\sindex[pkg]{#1}} % TBD: this must be extracted 
  % \newcommand{\tool}[1]{\texttt{#1}\sindex[tool]{#1}}% TBD: this must be extracted 
  % \newcommand{\param}[1]{\texttt{#1}\sindex[param]{#1}}% TBD: this must be extracted 
}

\title{The command \cmd{DocumentMetadata} and conformance with PDF/A and PDF/UE}
\author{Ernst Reissner (rei3ner@arcor.de)}
\date{\Today}


\begin{document}

\maketitle

\tableofcontents
\listoftables

\section{Introduction}\label{sec:intro}

This document is about using \cmd{DocumentMetadata}. 
This is the canonical way to force modern \LaTeX{} compilers to include metadata. 
Metadata in turn is the base for creating certain more strict flavors of PDF, 
namely PDF/A, PDF/X and PDF/UA\@. 
Whereas at time of this writing, PDF/X is not very much supported, 
this is different for the other standards. 

In general, one can specify the standard in \cmd{DocumentMetadata}, 
and the compiler tries to reach conformance, but there is no guarantee, 
even in absence of warnings. 
Thus, verification is mandatory and the canonical trustable and free verification tool 
with long term support is \texttt{verapdf} sponsored by the EU\@. 
Since this document also describes the circumstances under which conformance can be reached, 
we added a short description of \texttt{verapdf} in Section~\ref{sec:veraPdf}. 


%https://ctan.math.washington.edu/tex-archive/macros/latex-dev/required/latex-lab/documentmetadata-support-code.pdf
% parser for DocumentMetadata: 
% nested values only for key tagging-setup and debug but inside, no further nesting, 
% so regular expressions are still appropriate. 
% key may also comprise comments. 
% with only on level of nesting, we have w*<key>(={[~{}]})?(,\w*<key>(={[~{}]})?)*
% A bit restrictive, means no nested values, and no comments. 
%
%documentmetadata-support-doc.pdf 
% maybe in the manual a small section on pdfa: 
% neither hyperxmp nor pdfx is required. 
% just setting pdfstandard={A=2B}
% even colorscheme is not required. 
% Validation goes wrong for pdfversion={2.0} 
% but if nothing is given, it works out. 
% I think it does not work for pdflatex, but it works for lualatex 
% visual diff has differences but only in link colors: 
% internal links, links to bibliography and links to external three colors but now different. 
% we leave it as is, to indicate it must be PDF/A. 
% debug={xmp-export},
% one can do xmp-export=true or =\jobname.xmpi, nothing else, 
% because else the file cannot be erased. 
% TBD: what if standard is A-3x or even A=4, what if X- or UE=?
% debug={xmp-export} returns the same as pdfinfo -meta

% how shall we check pdfstandard? 
% we assume that neither pdfx nor hyperxmp are used. 
% only \DocumentMetadata. 
% The parser indicating latex main files 
% also finds \DocumentMetadata and the key pdfstandard. 
% Then https://ctan.math.illinois.edu/macros/latex/contrib/pdfmanagement-testphase/l3pdfmeta.pdf section 1.2:
%ocumentMetadata
% {
% %other options
% pdfstandard=A-2b,
% colorprofiles=
% {
% A = sRGB.icc, %or longer GTS_PDFA1 = sRGB.icc
% X = sRGB.icc,
% ISO_PDFE1=sRBG.icc
%
% another problem: TZ: latexmk always uses UTC do does ntlatex
% main application partially uses local time zone. 
% Then the according document is not reproducible if coming from another timezone 
% one would have to get from the pdf not only the time but also the timezone. 
% then one could create the new file with correct time and timezone. 
%

%\DocumentMetadata{lang=en,pdfversion=1.7,pdfstandard=A-2B}

% \DocumentMetadata{
%   %pdfversion={2.0}, %then not compliant
%   pdfstandard={A-2B},
%   %tagging={on}, %seems not available although documented
%   lang={en-US}, 
%   colorprofiles={A = sRGB.icc}
% }

\section{On \cmd{DocumentMetadata}}\label{sec:docMeta}


\section{On \texttt{verapdf}}\label{sec:veraPdf}% should remain section 3... else check dependencies. 

The list of return values given in Table~\ref{tab:exitCodesVeraPdf}. 
As can be seen, one code, 0, indicates that all checks were performed 
and all validation succeeded, one result means, the test could be performed and failed, 
and the other return values more or less indicate, 
that the proper tests could not be performed. 
Parsing errors means, it is no valid PDF file at all, so the proper tests could not be performed. 
Note that this application invoked \texttt{verapdf} with exactly one PDF file. 

\begin{table}[htb]
\centering
\begin{tabular}{lrl}
  \toprule
Enum                     & Exit code & Explanation \\
\midrule%
\midrule%
\Verb|VALID|                     &  0 & all files parsed and valid. \\    %VALID(0 & All files validated."), % also if option -o
\Verb|INVALID|                   &  1 & Invalid PDF/A file(s) found. \\
\Verb|BAD_PARAMS|                &  2 & Invalid command line parameters. \\
\Verb|OOM|                       &  3 & Out of Java heap space (memory). \\
\Verb|NO_FILES|                  &  4 & No files to process, also if no given files exist. \\% TBD: the latter must be added to docs of verapdf 
\Verb|IO_EXCEPTION|              &  6 & I/O Exception while processing. \\{}
--                               &  5 & Does not occur. \\
\Verb|FAILED_PARSING|            &  7 & Failed to parse one or more files. \\
\Verb|ENCRYPTED_FILES|           &  8 & Some PDFs encrypted. \\% TBD: what if --password is given but wrong? 
\Verb|VERAPDF_EXCEPTION|         &  9 & VeraPDF exception while processing. \\
\Verb|JAXB_EXCEPTION|            & 10 & Java XML marshalling exception \\
                                 &    &while processing result. \\
\Verb|FAILED_MULTIPROCESS_START| & 11 & Failed to start multiprocess. \\
\Verb|INTERRUPTED_EXCEPTION|     & 12 & Interrupted exception while processing. \\
\bottomrule%
\end{tabular}
\caption{\label{tab:exitCodesVeraPdf}The exit codes of \texttt{verapdf}}
\end{table}



\end{document}